\section{Native identification engines and the joint meaning of a counterfactual}
\label{exec19:cxxid-scope}
A lending analyst may ask whether an approval changes repayment, while a
servicing analyst asks whether a borrower who failed under the existing policy
would have succeeded under another action. These questions require different
information. A native implementation can make a symbolic calculation faster,
but its language does not turn an observational conditional distribution into
an intervention law or select a coupling between a person's potential outcomes.
The proposed product therefore keeps graph identification, numerical evaluation
and counterfactual model specification as separate contracts.

The selected aGrUM 3.2.0 revision is dated 6 October 2026. Twelve authenticated
implementation, wrapper and test originals extend the earlier dispatch review.
The build declarations require C++20 and CMake 3.18; the selected root license
statement is LGPL-3-or-later or MIT, subject to a separate component review.
These are historical source declarations, not a report that a native build,
Python installation or bank interface has passed. The R-backed C++ search core
of dosearch supplies a substantive alternative, while other native candidates
and partial-identification engines remain unresolved. The bounded comparison
cannot establish complete C++ market coverage.
\citep{exec19_cxx_identification_source,exec18_rid_dosearch}

\subsection{Inspecting recursion without mistaking it for a correctness proof}
\label{exec19:cxxid-recursion}
The native do-calculus source constructs backdoor and frontdoor expressions,
then provides a general recursive identification path. Its branches restrict
the graph to outcome ancestors, add irrelevant remaining variables to the
intervention, partition observed nodes connected through latent co-children,
and multiply or marginalize the resulting kernels. A single remaining full
confounded component triggers a hedge exception. The induced submodel keeps
observed identifiers and retains a latent descriptor only when it still has at
least two retained children. These observations identify the algorithmic
objects that a comparison must preserve; they are not a proof that every legal
input reaches the mathematically intended branch.
\citep{exec19_cxx_identification_source}

Two details deserve explicit qualification before reuse. The conditional-query
path builds one modified graph after removing incoming edges to intervention
nodes and outgoing edges from every conditioning node, then tests candidate
promotions within that graph. A comparison must establish the appropriate
multi-variable graph surgery and the conditional query after each promotion.
Separately, when an accumulated kernel is present, the first topological
factor is constructed as a joint-probability AST rather than the numerator
expression just calculated from that kernel. Whether that branch is reachable
with a distinct marginal, or is equivalent under the implementation's invariant,
requires a proof-aligned fixture. Neither observation alone warrants announcing
a software defect. Both are precise obligations that a root README or successful
simple adjustment example would miss.
\citep{exec19_cxx_identification_source}

The numerical layer has its own boundary. A formula's identified predicate
checks whether an expression root exists. Evaluation then requires an
observational Bayesian network, so a structure-only answer and an evaluable
answer are different states. Division of expression trees delegates to tensor
division. The selected files do not settle zero-denominator support handling,
finite arithmetic or an uncertainty statement. The Python convenience wrapper
retains the model on the returned tensor and maps an unidentified formula to a
\texttt{None} numerical answer; complete exception and lifetime conversions
still need direct inspection and execution. Preserve an expression with its
graph, variable domains and support conditions instead of returning an
unqualified floating-point success.
\citep{exec19_cxx_identification_source}

For example, let approval $A$ and repayment $Y$ be binary. In one authored
lending world, $U\sim\operatorname{Bernoulli}(1/2)$, $A=U$ and $Y=U$;
in another, $A\sim\operatorname{Bernoulli}(1/2)$ and $Y=A$. Both observed laws
put mass $1/2$ on each of $(A,Y)=(0,0)$ and $(1,1)$. Replacing the approval
equation gives an average intervention contrast of zero in the first world
and one in the second. Deterministic assignment in the first world also means
that the observed law does not provide all conditional treatment comparisons.
The example is a supplied structural construction, not a discovered bank graph.
Its purpose is to prevent an association or unsupported engine return from
being presented as identified repayment benefit.
\citep{exec19_cxx_identification_reference}

\subsection{Why individual root updates need a joint abduction contract}
\label{exec19:cxxid-abduction}
The selected counterfactual builder starts from the observational Bayesian
network, collects parentless nodes excluding interventions and declared latents,
conditions the original network on the observed profile, and fills each selected
root's conditional-probability table with its individual posterior. It then
calculates an intervention in the reconstructed model. The constructor used
for reconstruction does not copy the separate original latent descriptors or
spurious-arc declarations. Accordingly, graph identity also belongs in the
conformance test. More fundamentally, replacing several root marginals does
not in general retain their posterior joint dependence.
\citep{exec19_cxx_identification_source}

Consider a completely declared servicing example. Before an action, roots
$U$ and $V$ are independent fair bits; $A$ is an independent fair action.
Let success under action $a$ have one of two structural response functions:
\begin{align}
 Y_{\mathrm{stable}}(a)&=\mathbf 1\{U=V\},\nonumber\\
 Y_{\mathrm{switch}}(0)&=\mathbf 1\{U=V\},\qquad
 Y_{\mathrm{switch}}(1)=\mathbf 1\{U\ne V\}.
 \label{exec19:cxxid-coupling}
\end{align}
Every observed pair $(A,Y)$ has probability $1/4$ in both worlds; each world
also has $P\{Y(a)=1\}=1/2$ for both actions. Thus even the two marginal
intervention laws do not choose between these cross-world response functions.
For the profile $A=0,Y=1$, however, the roots must agree. The joint posterior
puts mass $1/2$ on $(U,V)=(0,0)$ and $1/2$ on $(1,1)$, while each separate
root posterior remains a fair bit. Consequently,
\begin{align}
 P\{Y_{\mathrm{stable}}(1)=1\mid A=0,Y=1\}&=1,\nonumber\\
 P\{Y_{\mathrm{switch}}(1)=1\mid A=0,Y=1\}&=0,\nonumber\\
 \sum_{u,v}\mathbf 1\{u=v\}P(u\mid A=0,Y=1)P(v\mid A=0,Y=1)&=\tfrac12.
 \label{exec19:cxxid-joint}
\end{align}
The last expression is the result of imposing separate posterior root
marginals; it is not the first world's joint conditional probability. In the
switch world the same product reconstruction also gives $1/2$ for inequality.
The exact finite enumeration tests this semantic distinction without executing
an upstream package. A native run is still required to establish the behavior
of the installed representation, graph and wrapper on this fixture.
\citep{exec19_cxx_identification_reference}

If the declared payoff is USD 100 for success less a USD 20 action cost, the
three respective success probabilities imply USD 80, USD $-20$ and USD 30.
These are conditional values for the specified profile and action, not an
incremental return against a measured current workflow. They are assumed
monetary consequences of a semantic choice. A product cannot substitute the
middle probability $1/2$ because it looks numerically stable: it must preserve
the joint abduction distribution, state a sufficient factorization condition,
return a justified set over admissible couplings, or refuse the requested
individual answer. Additional marginal intervention observations alone do not
resolve this constructed coupling ambiguity.
\citep{exec19_cxx_identification_reference}

\subsection{A reusable boundary and a concrete qualification route}
\label{exec19:cxxid-contract}
The inspected native tests are useful starting specifications. They include
manual backdoor/frontdoor expression comparisons, selected direct-do-calculus
numerical comparisons, a hedge evaluation exception, ID/IDC smoke tests and a
salary counterfactual fixture. The selected conditional smoke tests verify that
an expression builds and evaluates; they do not independently establish its
multi-variable causal meaning. A parity check against the same builder also
cannot replace a separate structural oracle. None of these upstream tests was
run in this source review, and the coupled-root example is a new authored
reference, not a claim about existing upstream coverage.
\citep{exec19_cxx_identification_source,exec19_cxx_identification_reference}

For dosearch, the already pinned R/C++ source makes a different tradeoff. Its
native entry accepts directed and bidirected graph information, a collection
of known distributions, a target and search options through Rcpp. The loop can
return when its time budget expires; the returned identifiable flag depends
on whether a target distribution was found. The inspected result does not
separately certify search exhaustion. Therefore preserve a discovered formula
and derivation when available, but retain unknown or time-limited status when
search failure has not been proved to imply nonidentification in the stated
problem class. Start reuse through the supported R boundary; a stand-alone C++
port would additionally require extraction of the Rcpp interface, environment,
representation and rights obligations. A shared implementation language is
insufficient evidence that the two engines deliver an interchangeable service.
\citep{exec18_rid_dosearch}

The proposed conformance program first fixes the query, graph class, observed
or experimental input law, domains, support and information clock. It then
compares direct and shortcut identification on independently enumerated small
models, tests multi-variable conditioning and nested kernels, exercises
nonidentification separately from unsupported methods, and tests numeric
availability and support-zero behavior. Counterfactual tests additionally
retain structural response functions and joint posterior dependence; a graph
with ordinary conditional tables is not silently promoted to a unique
cross-world model. Finally, the same signed query and typed result cross the
native, Python and process boundaries, with ownership and exceptions checked.
Retain aGrUM and dosearch as conditional comparators and defer executable
admission until those trials pass in their declared regimes. Other candidates,
original proof alignment, full component rights and general partial
identification remain open; this bounded source review does not close them.
\citep{exec19_cxx_identification_source,exec18_rid_dosearch,exec19_cxx_identification_reference}
